\chapter{الخاتمة}
\label{ch:conclusion}


قدم هذا المشروع دراسة هندسية متكاملة في مجال \textbf{التحكم بالحركة (\textenglish{Motion Control})}، متخذاً من منظومة "المقص الطائر" (\textenglish{Flying Shear}) تطبيقاً عملياً متقدماً لاختبار النظريات الديناميكية في خطوط الإنتاج المستمرة. انطلق البحث من التأسيس النظري لمعمارية التحكم الحركي وحلقات التحكم المتداخلة (\textenglish{Cascaded Loops})، وصولاً إلى دراسة التحدي الأكبر في هذا المجال: تحقيق المزامنة الكينماتيكية (\textenglish{Kinematic Synchronization}) بين محورين مستقلين (المحور القائد للسرعة، والمحور التابع للتموضع).

من خلال هذا العمل، تم إثبات فعالية تحويل المعادلات الديناميكية إلى خوارزميات منطقية قابلة للتنفيذ. حيث تم توظيف "التحكم النبضي" عالي التردد جنباً إلى جنب مع العدادات السريعة (32-بت) في بيئة متحكمات \textenglish{Delta PLC} لفرض مسار حركة دقيق (\textenglish{Trajectory Generation}). كما أثبتت خوارزمية المخطط السلمي (\textenglish{Ladder Logic}) نجاحها في معالجة المتغيرات الحركية اللحظية؛ بدءاً من اتخاذ قرار بدء الحركة، وتطبيق التصفير الذاتي للمشفرات لتجنب تراكم أخطاء التموضع (\textenglish{Positioning Errors})، وانتهاءً بضمان الكبح الديناميكي الآمن عبر شبكات حماية نهايات الشوط.

يُمثل هذا المزج بين التصميم المنطقي والاختبار المعمق عبر بيئة المحاكاة، دليلاً عملياً على إمكانية تطوير حلول تحكم حركي موثوقة وعالية الأداء، مما يعكس نجاح المشروع في سد الفجوة بين نظريات الأتمتة المتقدمة وتطبيقاتها الصناعية الدقيقة.


\vspace{2cm}
\textbf{لمشاهدة منظومة الآلة عملياً:}
\begin{center}
\LR{https://youtu.be/uw45qdJEOuA}
\end{center}
